\documentclass[10pt,a4paper]{amsart}
\usepackage[margin=19mm]{geometry}
\usepackage{amsmath,amssymb,mathtools}
\usepackage[scr=rsfs,cal=euler]{mathalfa}
\usepackage{xfrac}
\usepackage[unicode,hidelinks,bookmarks=false]{hyperref}
\newcommand{\ie}{i.e.\ }
\newcommand{\cf}{cf.\ }
\newcommand{\resp}{respectively\ }
\newcommand{\Subsection}{\subsection}
\providecommand{\nobreakdash}{\nobreak\hbox{-}}
\setlength{\emergencystretch}{2em}
\multlinegap0pt
\usepackage{amssymb}
\usepackage{xcolor}
\newtheorem{lemma}{Lemma}
\newtheorem{theorem}{Theorem}
\newcommand{\Class}{\mathcal{C}}
\newcommand{\Eadapt}{E_{\mathrm{adapt}}}
\newcommand{\Erasure}{c}
\newcommand{\Fover}{F}
\newcommand{\Hw}{\mathcal{H}}
\newcommand{\MutInfo}{I_{\mathrm{alg}}}
\newcommand{\Nop}{N}
\newcommand{\Sys}{\Sigma}
\newcommand{\Task}{\mathcal{T}}
\newcommand{\descop}{\mathrm{desc}}
\title{Watts-per-Intelligence Part II:\\Algorithmic Catalysis}
\author{Elija Perrier}
\date{Source: arXiv:2604.20897v2 (CC BY 4.0)}
\begin{document}
\maketitle

\noindent\emph{Source and licence.} Watts-per-Intelligence Part II:\\Algorithmic Catalysis. Version arXiv:2604.20897v2 (CC BY 4.0), distributed by arXiv under the Creative Commons Attribution 4.0 International licence, http://creativecommons.org/licenses/by/4.0/, as recorded in the arXiv licence metadata for this exact version and shown on the abstract page https://arxiv.org/abs/2604.20897v2. Full licence notice: \texttt{licenses/arxiv-2604.20897-CC-BY-4.0.txt}.\par\smallskip

\noindent\textit{Editorial scope.} Counted unit: the article's theorem thm:composition (source0.tex lines 671--696). The article's complete original proof of it is delivered with it (source0.tex lines 697--717, one \texttt{proof} environment of 21 lines). No mathematics is written, repaired or re-derived: the statement and the proof are the article's own lines, in the article's own notation, and no retained line is substituted or re-flowed.  Retained uncounted as the source-local context the proof needs: lines 293--299; lines 560--572.  Nothing was omitted from any retained span, so the package carries no omission record.  No retained line contains a citation, so the wrapper carries no bibliography and no bibliography entry.  No retained line contains a figure environment, an image-inclusion command or drawing code, so no image is delivered and the package declares no asset dependency.  Editorial formatting only, in addition: an AMS \texttt{amsart} preamble that restates the article's own declarations and macros used by the retained lines, namely the environments \texttt{lemma}, \texttt{theorem} on separate counters, together with the standard packages those lines need.  The retained environments print in this document as Lemma 1, Theorem 1 in order of appearance, whereas the article prints the same items with the numbering of its own full document.  The complete native source of the article as distributed in the arXiv e-print for this exact version is bundled unchanged as \texttt{sources/198983/source0.tex}.
\begin{lemma}[Substrate monotonicity]
\label{lem:monotone}
If $\Hw'$ refines $\Hw$, then
$\MutInfo(\descop(\Hw'):\sigma(\Class))\ge
\MutInfo(\descop(\Hw):\sigma(\Class))-c_U$, and in particular
$\eta'\ge\eta-c_U$.
\end{lemma}

\begin{theorem}[Thermodynamic cost of structural information]
\label{thm:adaptation-cost}
The adaptation energy $\Eadapt$ on $\Hw_{\mathrm{adapt,cat}}$ of any system
$\Sys_{\mathrm{cat}}$ with post-adaptation substrate mutual
information $\mu$ and adaptation input $\mathcal{D}$ satisfies
\begin{equation}
\label{eq:adaptation-floor}
\Eadapt\;\ge\;
\Fover(\Hw_{\mathrm{adapt,cat}})\,\Erasure\,
\bigl[\mu\,-\,\MutInfo(\mathcal{D}:\sigma(\Class))\,-\,c_U\bigr]_+,
\end{equation}
where $[x]_+:=\max(x,0)$.
\end{theorem}

\begin{theorem}[Composition of catalysts]
\label{thm:composition}
Let $\Sys_0\to\Sys_1\to\Sys_2$ be a chain of algorithmic catalysts
on a common task class $\Class$, with stagewise class-specific
speed-ups $\Gamma_1,\Gamma_2$ and structural-information parameters
$\eta_1,\eta_2$, and assume that $\descop(\Hw_{\mathrm{exec},2})$
refines $\descop(\Hw_{\mathrm{exec},1})$ in the sense of
Lemma~\ref{lem:monotone}. Then the composite system
$\Sys_0\to\Sys_2$ satisfies
\begin{equation}
\label{eq:composition}
\Gamma_{1\circ 2}\;\ge\;\Gamma_1\cdot\Gamma_2,
\qquad
\eta_{1\circ 2}\;\ge\;\max\{\eta_1,\eta_2\}-c_U,
\end{equation}
with the stronger additive bound
$\eta_{1\circ 2}\ge\eta_1+\eta_2-c_U$ holding when the two stages
encode algorithmically independent aspects of $\sigma(\Class)$, in
the sense that $\MutInfo(\descop(\Hw_{\mathrm{exec},1}):
\descop(\Hw_{\mathrm{exec},2})\mid\sigma(\Class))=O(1)$. The
adaptation energy lower bound for the composite is
$\Eadapt^{1\circ 2}\ge\Fover(\Hw_{\mathrm{adapt},2})\Erasure
[\eta_{1\circ 2}-\MutInfo(\mathcal{D}_{\mathrm{tot}}:\sigma(\Class))
-c_U]_+$, where $\mathcal{D}_{\mathrm{tot}}$ is the concatenated
adaptation input of the two stages.
\end{theorem}

\begin{proof}
Multiplicativity of $\Gamma$ follows from the definitional
identity
$\Nop(\Sys_0,\Task)/\Nop(\Sys_2,\Task)
=(\Nop(\Sys_0,\Task)/\Nop(\Sys_1,\Task))\cdot
(\Nop(\Sys_1,\Task)/\Nop(\Sys_2,\Task))\ge\Gamma_1\Gamma_2$ at
matched intelligence. The refinement hypothesis and
Lemma~\ref{lem:monotone} yield $\eta_{1\circ 2}\ge\eta_1-c_U$ and
$\eta_{1\circ 2}\ge\eta_2-c_U$, hence
$\eta_{1\circ 2}\ge\max\{\eta_1,\eta_2\}-c_U$. For the independent
case, algorithmic independence of
$\descop(\Hw_{\mathrm{exec},1})$ and $\descop(\Hw_{\mathrm{exec},2})$
conditional on $\sigma(\Class)$ implies, up to $c_U$, that
$\MutInfo(\descop(\Hw_{\mathrm{exec},2}):\sigma(\Class))\ge
\MutInfo(\descop(\Hw_{\mathrm{exec},1}):\sigma(\Class))
+\MutInfo(\descop(\Hw_{\mathrm{exec},2}):\sigma(\Class)\mid
\descop(\Hw_{\mathrm{exec},1}))$ by the algorithmic chain rule,
and each term on the right lower-bounds the corresponding
$\eta_i$ up to $c_U$. The adaptation-energy bound follows from
Theorem~\ref{thm:adaptation-cost} applied to the composite.
\end{proof}

\end{document}
